% from ega4-errata-addenda-list3-fr.tex
\oldpage[IV]{354}
Soit alors $X'$ le normalisé de $X$ (II, 6.3.8); si $X_i$
($1\leq i\leq n$) sont les sous-préschémas réduits de $X$ ayant pour
espaces sous-jacents les composantes irréductibles de $X$, $X'$ est
somme des normalisés $X'_i$ des préschémas intègres $X_i$; nous
désignerons par $x_i$ (resp. $x'_i$) le point générique de $X_i$
(resp. $X'_i$). Le morphisme canonique $p:X'\to X$ est fini puisque
$X$ est universellement japonais et on a
$p^{-1}(x_i)=\{x'_i\}$, l'homomorphisme $k(x_i)\to k(x'_i)$ étant
bijectif.

Soit $f=(\psi,\theta)$; on sait déjà
{\hyperref[canon-index.ega4.c04.statement.10-4-11]{(10.4.11)}} que $\psi$ est une
bijection continue de l'espace $X$ sur lui-même; si l'on pose
$y_i=\psi(x_i)$, $y_i$ est aussi un point maximal de $X$, car si $y'_i$
était une générisation de $y_i$ distincte de $y_i$, $x_i$ serait une
générisation $\neq x_i$ de $\psi^{-1}(y'_i)$. Puisque l'ensemble $M$
des points maximaux de $X$ est fini, la restriction $\psi|M$ est une
bijection de $M$ sur lui-même, et on peut écrire
$\psi(x_i)=x_{\sigma(i)}$, où $\sigma$ est une permutation de
$[1,n]$. Alors, pour tout $i$, la restriction $f|X_i$ de $f$ se
factorise en
\[
 X_i\xrightarrow{f_i}X_{\sigma(i)}\longrightarrow X
\]
(I, 5.2.2), où la seconde flèche est l'injection canonique. Comme les
$X_i$ sont intègres, il existe (II, 6.3.9) un unique morphisme
$f'_i:X'_i\to X'_{\sigma(i)}$ rendant commutatif le diagramme
\[
\begin{array}{ccc}
X'_i&\xrightarrow{f'_i}&X'_{\sigma(i)}\\
\downarrow p_i&&\downarrow p_{\sigma(i)}\\
X_i&\xrightarrow{f_i}&X_{\sigma(i)}
\end{array}
\]
et si $f'$ coïncide avec $f'_i$ dans chacun des $X'_i$, on a donc un
diagramme commutatif
\phantomsection\EGAReaderTarget{ega4.c04.sec.10-4-12-1}{IV.10.4.12.1}
\begin{equation}
\label{IV.10.4.12.1.errata}
\begin{array}{ccc}
X'&\xrightarrow{f'}&X'\\
\downarrow p&&\downarrow p\\
X&\xrightarrow{f}&X.
\end{array}
\tag{10.4.12.1}
\end{equation}

D'autre part, comme $f$ est séparé, il en est de même de
$f\circ p=p\circ f'$, donc $f'$ est séparé (I, 5.5.1); enfin, puisque
$p$ est fini et $f$ bijectif, il résulte de
{\hyperlink{ega4.c04.sec.10-4-12-1}{(10.4.12.1)}} que $f'$ est
quasi-fini. Puisque les $X'_i$ sont intègres et normaux, il résulte de
{\hyperref[canon-index.ega4.c04.cor.8-12-11]{(8.12.11)}} que $f'_i$ se factorise en
\[
 X'_i\xrightarrow{f''_i}X''_{\sigma(i)}
 \xrightarrow{u_{\sigma(i)}}X'_{\sigma(i)},
\]
où $X''_{\sigma(i)}$ est la fermeture intégrale de
$X'_{\sigma(i)}$ relative à l'extension $k(x'_i)$ de
$k(x'_{\sigma(i)})$, et $f''_i$ une immersion ouverte. Mais comme
$k(x'_i)$ est une extension radicielle de $k(x'_{\sigma(i)})$,
$u_{\sigma(i)}$ est un morphisme radiciel (Bourbaki,
\emph{Alg. comm.}, chap.~V, \S2, n\textsuperscript{o}~3, lemme 4);
donc $f'$ est radiciel, et il résulte alors de
{\hyperref[canon-index.ega4.c04.statement.10-4-11]{(10.4.11)}} que $f'$ est
surjectif; on en conclut (puisque les $u_{\sigma(i)}$ sont injectifs)
que $f''_i$ est bijectif, donc un isomorphisme; on en déduit enfin que
le morphisme $f'$ est fini (puisqu'il en est ainsi de $u_{\sigma(i)}$,
$X'$ étant universellement japonais
{\hyperref[canon-index.ega4.c04.statement.7-8-3]{(7.8.3)}}). Cela étant, comme
$p\circ f'=f\circ p$ est alors fini et que $p$ est surjectif et $f$
séparé, $f$ est propre (II, 5.4.3) et quasi-fini, donc fini
{\hyperref[canon-index.ega4.c04.thm.8-11-1]{(8.11.1)}}.

